%! Factoring Special Forms — Unit 1, Lesson 1.5 — Group Activity
\documentclass[10pt]{article}
\usepackage{atda-article}
\usepackage{atda-boxes}
\newcommand{\TallMath}[1]{$\displaystyle #1\rule[-1.4em]{0pt}{3.2em}$}
\setlength{\workrowsep}{6pt}

\begin{document}

\pageheader{Unit 1, Lesson 1.5}{Group Activity}
% NAMESTRIP: no \namedateperiod here. The student writes their name once, on the
% cover this component is stapled behind. See references/conventions.md.

\vspace{0.15in}

\begin{scenariobox}[The Scenario]{royal}
\small
The drama club is ordering a \textbf{rectangular plywood platform} for the spring show. As usual,
the shop teacher hands over the \emph{area} and nothing else --- and the club has to order edge trim
before the build starts. Same job as last week, one difference: today the polynomial usually has
only \textbf{two} terms, so there is nothing to group and no middle term to split.

Work with your group. \textbf{Everyone starts at Tier R.} When your whole group can do Tier R
without help, move on.
\end{scenariobox}

\vspace{0.10in}

\boxguard[16]
\begin{tcolorbox}[colback=white, colframe=black!40, title=\textbf{Tier R --- Remediate},
    fonttitle=\bfseries, arc=2mm, left=3mm, right=3mm, top=2mm, bottom=2mm, breakable]
\small
Five fluency moves --- one for each pattern, and one more. Name the pattern before you factor, and
check every answer by multiplying back.

\begin{enumerate}[leftmargin=1.6em, itemsep=4pt, topsep=4pt]
  \item Factor: $x^2-144$
\begin{work}
  x^2-144 &= x^2-12^2 \\
          &= (x+12)(x-12)
\end{work}

  \item Factor: $x^2+14x+49$
\begin{work}
  2(x)(7) &= 14x \quad\checkmark \\
  x^2+14x+49 &= (x+7)^2
\end{work}

  \item Factor: $x^2-10x+25$
\begin{work}
  2(x)(5) &= 10x \quad\checkmark \\
  x^2-10x+25 &= (x-5)^2
\end{work}

  \item Factor: $9x^2-4$
\begin{work}
  9x^2-4 &= (3x)^2-2^2 \\
         &= (3x+2)(3x-2)
\end{work}

  \item Factor: $x^3+27$ \quad (two terms, both perfect cubes --- use SOAP)
\begin{work}
  x^3+27 &= x^3+3^3 \\
         &= (x+3)\left(x^2-3x+9\right)
\end{work}
\end{enumerate}
\end{tcolorbox}

\vspace{0.10in}

\boxguard[16]
\begin{tcolorbox}[colback=white, colframe=black!40,
    title=\textbf{Tier A --- Approaching Proficiency},
    fonttitle=\bfseries, arc=2mm, left=3mm, right=3mm, top=2mm, bottom=2mm, breakable]
\small
The approved platform has area
\[ A(x) = 9x^2-25 \quad \text{square feet.} \]

\begin{enumerate}[leftmargin=1.6em, itemsep=4pt, topsep=4pt]
  \item Factor $A(x)$ to recover the platform's two side lengths.
\begin{work}
  9x^2-25 &= (3x)^2-5^2 \\
          &= (3x+5)(3x-5)
\end{work}

  \item The two sides are not equal, but they are the same distance from $3x$ feet --- one longer,
        one shorter. How far, and how does the factored form tell you?

\writeline

  \item The club builds at $x=4$. Find both side lengths, the area, and the edge trim the club must
        order --- then check your area against the original polynomial.
\begin{work}
  3x+5 &= 17 \text{ ft} \qquad 3x-5 = 7 \text{ ft} \\
  17 \cdot 7 &= 119 \text{ sq ft} \\
  9(16)-25 &= 119 \text{ sq ft} \quad\checkmark \\
  \text{trim} &= 2(17)+2(7) = 48 \text{ ft}
\end{work}

  \item A second platform is offered with area $A(x)=16x^2-81$. Factor it and give both dimensions.
\begin{work}
  16x^2-81 &= (4x)^2-9^2 \\
           &= (4x+9)(4x-9)
\end{work}

  \item A third is quoted as $A(x)=5x^2-45$. Factor it \textbf{completely} --- look at what the two
        terms share before you reach for a pattern.
\begin{work}
  5x^2-45 &= 5\left(x^2-9\right) \\
          &= 5(x+3)(x-3)
\end{work}
\end{enumerate}
\end{tcolorbox}

\vspace{0.10in}

\boxguard[30]
\begin{tcolorbox}[colback=white, colframe=black!40, title=\textbf{Tier E --- Extension},
    fonttitle=\bfseries, arc=2mm, left=3mm, right=3mm, top=2mm, bottom=2mm, breakable]
\small

\begin{enumerate}[leftmargin=1.6em, itemsep=4pt, topsep=4pt]
  \item \textbf{Critique the work.} A student turned in these two answers. Neither is acceptable.
        Write the correct answer for each, then name the mistake.
        \[ x^2+64 = (x+8)(x-8) \qquad\text{and}\qquad x^2-6x+9 = (x-3)(x+3) \]
\begin{work}
  (x+8)(x-8) &= x^2-64 \ne x^2+64 \\
  x^2+64 &\text{ is a sum of squares, so it is prime} \\
  x^2-6x+9 &= (x-3)^2
\end{work}
\writelines{2}

  \item \textbf{One pattern, twice.} A backdrop panel is quoted as $A(x)=x^4-81$ square feet. Factor
        it \textbf{completely}, and say which factor stops and why.
\begin{work}
  x^4-81 &= \left(x^2\right)^2-9^2 \\
         &= \left(x^2+9\right)\left(x^2-9\right) \\
         &= \left(x^2+9\right)(x+3)(x-3)
\end{work}

  \item \textbf{Verify the identity.} The difference-of-cubes rule claims that
        $a^3-b^3 = (a-b)\left(a^2+ab+b^2\right)$ for \emph{every} pair of numbers. Multiply the
        right-hand side out to check it, then explain why doing this once with letters settles it
        for every pair of numbers at once.
\begin{work}
  (a-b)\left(a^2+ab+b^2\right) &= a^3+a^2b+ab^2-a^2b-ab^2-b^3 \\
                               &= a^3-b^3
\end{work}
\writelines{2}
\end{enumerate}
\end{tcolorbox}

\end{document}
